In order to answer to the previous research question, we distinguish six steps in this 
master thesis. These steps can be observed in the Figure \ref{fig:roadmap}.

\subsection*{P\&ID and Preliminary Thermodynamic cycle design}

In this first part, we will numerically modeled the thermodynamical design. Each student 
will work simultaneously on different blocks, determined in advance. The analysis of the 
cycle will include an energy and exergy analysis to ensure optimal performances. The numerical 
model will be writen in Python and the module Coolprop will be used to access numerical
thermodynamics tables. 

\subsection*{Optimization of the Thermodynamic Design}

Assembling all the different blocks will be the purpose of the second part. Beside 
connecting the blocks, an optimization process will be developped to determine the 
optimal parameters in nominal conditions. In addition, more physical aspects will be 
added to the model such as head losses computations and heat transfer modeling. 

\subsection*{Sizing of the Main Components}

Once the thermodynamic design is fixed, the main components as heat exchangers, expansion
valves and water tanks will be chosen. 

\subsection*{Establishment of a First Control Strategy}

One main advantage of this machine will be its capacity of simulating different conditions
thanks to the dual evaporator. A control logic must therefore be implemeted. This part
is critical because we'll have to think about the nature and the location of the sensors. 
\subsection*{CAD and Operational Model}

In this part, the other student will be in charge of the design of the machine. It will 
contain material/components selection, heat exchangers and insulation design and the 
elaboration of precise technical draw. That will be our task to correct the final cycle design 
(head losses, heat exchange, \dots). 

\subsection*{Prototype Construction and Instrumentation} 

The current part will be the construction of the prototype.

\subsection*{Safety Validation of the Prototype}

Before starting the tests, the safety of the machine will be verified by the CREDEM platform.

\subsection*{Control Logic and Performance Validation}

Finally, the last part will be dedicated to the evaluation of the machine and its control logic. 
Several tests will take 
place to ensure the good operation of the machine and to collect data in order to compare them 
to theoritical values. 


